% ============================================================
% v2 — Capitolo 1: Introduzione
% Blocco I di IV (intro -> progetto -> evoluzioni -> conclusione)
% Budget: ~4 pagine. Include i fondamenti essenziali, ridotti a
% ciò che serve effettivamente nei capitoli successivi.
% ============================================================
\chapter{Introduzione}
\label{ch:intro}

\section{Motivazione e domanda di ricerca}
\label{sec:motivazione}

L'apprendimento per rinforzo (\emph{reinforcement learning}, RL)
addestra politiche di controllo per tentativi ed errori, e in qualunque
dominio non banale il progettista incontra una decisione di scheduling
prima ancora del primo passo di gradiente: quando i task da padroneggiare
coprono un intervallo di difficoltà, \emph{in quale ordine vanno
presentati}? Due risposte sono naturali e opposte. Il \emph{curriculum
learning} li presenta dal facile al difficile, sulla falsariga della
pedagogia umana: si costruisce competenza dove il successo è
raggiungibile, poi la si estende. L'\emph{addestramento misto}
(mixed-batch) li presenta tutti fin dall'inizio, ed è l'analogo diretto
dell'addestramento su un dataset mescolato, default indiscusso
dell'apprendimento supervisionato.

Le due strategie consumano lo stesso budget di calcolo: scegliere tra
loro è, in linea di principio, gratuito. Se la scelta cambia ciò che la
policy appresa fa, allora è tra le leve più economiche a disposizione di
chi progetta un sistema di RL --- ed è anche tra le meno caratterizzate
nei contesti \emph{multi-agente}, dove quasi nessuna delle assunzioni
della letteratura sui curricula (apprendista singolo, reward sparse,
ambiente stazionario) continua a valere.

La guida autonoma urbana è un banco di prova esigente per porre la
domanda: offre un asse di difficoltà naturale e continuo (la lunghezza
del percorso), agenti eterogenei con vincoli fisici diversi, una
tensione strutturale tra prestazione e sicurezza --- le policy più
veloci completano più percorsi \emph{e} urtano più cose --- e un test
out-of-distribution canonico, guidare su una città mai vista in
addestramento. È inoltre intrinsecamente multi-agente: le policy non si
limitano a condividere un mondo, ma plasmano l'ambiente l'una
dell'altra mentre apprendono.

\begin{quote}
  \emph{Il curriculum learning da easy a medium a hard produce un
  comportamento di guida multi-agente misurabilmente diverso rispetto
  all'addestramento mixed-batch, a parità di budget di addestramento?}
\end{quote}

La domanda riguarda deliberatamente il \emph{comportamento} e non una
singola classifica di tassi di successo: il confronto è valutato lungo
quattro assi --- efficienza campionaria iniziale, prestazione finale con
la sua composizione dei fallimenti, dinamica entro la run e
generalizzazione a una mappa nuova --- sempre separando veicoli e
pedoni. Tre ipotesi falsificabili strutturano l'analisi:

\begin{itemize}
  \item \textbf{H1 (curriculum).} L'esposizione graduale produce una
        migliore efficienza campionaria iniziale, perché i task facili
        offrono successi raggiungibili da cui partire; con il rischio
        noto di plateau od oblio sui task difficili introdotti tardi.
  \item \textbf{H2 (misto).} L'esposizione uniforme copre l'intera
        distribuzione dei task dal passo zero, al costo di un segnale
        iniziale più rumoroso; il suo payoff atteso è la robustezza.
  \item \textbf{H3 (interazione, secondaria).} L'entità --- forse la
        direzione --- dell'effetto di regime può dipendere da come viene
        stimato il valore: nelle architetture ad addestramento
        centralizzato è il \emph{critic} a digerire la struttura
        multi-agente, e sostituirne l'encoder può cambiare ciò che
        l'ordinamento produce.
\end{itemize}

\section{Fondamenti essenziali}
\label{sec:fondamenti}

Questa sezione introduce solo i concetti effettivamente usati nel
seguito, ciascuno alla profondità alla quale serve.

\paragraph{RL, valore e vantaggio.}
Il RL formalizza la decisione sequenziale come processo decisionale di
Markov $(\mathcal{S}, \mathcal{A}, P, r, \gamma)$~\cite{sutton2018rl}:
un agente agisce tramite una policy $\pi(a \mid s)$ e massimizza il
ritorno atteso scontato
$J(\pi) = \mathbb{E}_{\pi}[\sum_{t \geq 0} \gamma^{t} r(s_t,a_t)]$. Lo
sconto fissa un orizzonte decisionale effettivo di $\approx 1/(1-\gamma)$
passi: non è una costante cosmetica, perché la piattaforma del
\cref{ch:piattaforma} usa $\gamma = 0.997$, cioè $\approx 333$ passi, su
episodi lunghi fino a mille passi. Il \emph{vantaggio}
$A^{\pi}(s,a) = Q^{\pi}(s,a) - V^{\pi}(s)$ misura quanto un'azione sia
migliore della media della policy in quello stato, e il teorema del
gradiente di policy fornisce
$\nabla_{\theta} J = \mathbb{E}_{\pi_\theta}[\nabla_{\theta}\log\pi_{\theta}(a_t\mid s_t)\,A^{\pi}(s_t,a_t)]$.
Poiché stimare $A$ con ritorni Monte Carlo grezzi produce alta varianza,
si sottrae una baseline appresa: i metodi \emph{actor--critic}
addestrano perciò l'actor $\pi_\theta$ e il critic $V_\phi$ --- e questa
divisione del lavoro è l'aggancio all'architettura multi-agente, perché
il critic è un dispositivo del solo addestramento e può quindi ricevere
informazione che l'actor non vede mai.

\paragraph{PPO, GAE ed entropia.}
Gli aggiornamenti di gradiente ``vanilla'' sono fragili, perché un
singolo passo troppo grande può far collassare una policy funzionante.
PPO~\cite{schulman2017ppo} lo evita con un obiettivo surrogato clippato:
detto $\rho_t(\theta) = \pi_{\theta}(a_t\mid s_t)/\pi_{\theta_{\text{old}}}(a_t\mid s_t)$
il rapporto di importanza,
\begin{equation}
  L^{\mathrm{CLIP}}(\theta) =
  \mathbb{E}_t\!\left[\min\bigl(\rho_t A_t,\;
  \operatorname{clip}(\rho_t,\, 1-\varepsilon,\, 1+\varepsilon)\, A_t\bigr)\right],
\end{equation}
il che rimuove ogni incentivo a spostare il rapporto oltre
$[1-\varepsilon, 1+\varepsilon]$ e permette più epoche di gradiente su
ogni batch senza aggiornamenti distruttivi. I vantaggi $A_t$ provengono
da GAE~\cite{schulman2016gae}, che con i residui temporali
$\delta_t = r_t + \gamma V(s_{t+1}) - V(s_t)$ pone
$A_t = \sum_{l \geq 0} (\gamma\lambda)^{l}\delta_{t+l}$, dove $\lambda$
scambia bias e varianza tra il TD a un passo e i ritorni Monte Carlo, e
la cui qualità dipende direttamente da quella del critic. Un bonus di
entropia sostiene infine l'esplorazione penalizzando il collasso
prematuro della policy.

\paragraph{MARL, CTDE e MAPPO.}
Con $N$ agenti in un unico ambiente l'MDP si generalizza a un
\emph{gioco di Markov}, e sorgono due difficoltà strutturali. La prima è
la \emph{non-stazionarietà}: dal punto di vista di ciascun agente gli
altri sono parte dell'ambiente e stanno apprendendo, quindi la dinamica
effettiva cambia nel tempo e l'esperienza raccolta sotto vecchie policy
dei co-giocatori diventa fuorviante~\cite{lowe2017maddpg}. La seconda è
l'\emph{attribuzione del merito}: il ritorno di ogni agente dipende
dalle azioni di tutti. La risoluzione ormai standard sfrutta la
divisione actor--critic: si addestrano i critic con accesso a
informazione globale, mantenendo gli actor condizionati sulle sole
osservazioni locali~\cite{lowe2017maddpg}. Il critic assorbe la
non-stazionarietà, mentre l'esecuzione resta pienamente decentralizzata
perché i critic vengono scartati dopo l'addestramento: è il paradigma
\emph{centralized training, decentralized execution} (CTDE), usato in
tutta questa tesi. MAPPO~\cite{yu2022surprising} lo istanzia con actor
PPO e funzione di valore centralizzata, e Yu et al.\ mostrano che questa
combinazione minimale eguaglia o supera i metodi off-policy sui
benchmark cooperativi standard purché i dettagli implementativi siano
corretti --- fra i quali proprio il design dell'input del critic. Nel
critic più semplice quell'input è la concatenazione piatta delle
osservazioni di tutti gli agenti; le alternative strutturate le trattano
come nodi di un grafo, ed è il confronto tra queste due famiglie a
costituire l'asse secondario di questa tesi (\cref{sec:critic}).

\paragraph{Curriculum learning e la sua condizione di controllo.}
L'idea che l'apprendimento migliori quando gli esempi sono presentati in
un ordine sensato precede il deep learning: gli esperimenti
\emph{starting small} di Elman mostrarono reti ricorrenti capaci di
apprendere una grammatica complessa solo quando la complessità
dell'input cresceva incrementalmente~\cite{elman1993starting}, e Bengio
et al.~\cite{bengio2009curriculum} diedero poi alla strategia il nome di
curriculum learning, inquadrandola come metodo di continuazione in cui
gli esempi facili definiscono una versione più liscia dell'obiettivo.
Nel RL il progettista controlla non solo
l'ordine dei dati ma la distribuzione dei \emph{task}, il che rende i
curricula al tempo stesso più potenti e più difficili da progettare;
nella tassonomia della survey di Narvekar et
al.~\cite{narvekar2020curriculum} il meccanismo della \cref{sec:regimi}
è uno schema di \emph{sequenziamento} basato sulla prestazione su un
insieme fisso di task. La motivazione riguarda in larga parte
l'esplorazione: sui task difficili una policy debole può non
sperimentare praticamente mai il successo, e quindi non ricevere segnale
di apprendimento. Tre modi di fallimento noti hanno plasmato il design
del \cref{ch:piattaforma} --- l'\emph{oblio} della competenza precedente
sotto distribuzione spostata~\cite{french1999catastrophic}, lo
\emph{stallo} su un gate mai soddisfatto e la \emph{mis-calibrazione}
del gate stesso --- e sono esattamente i rischi che H1 mette alla prova.
Ogni affermazione sull'ordinamento richiede però una baseline in cui
esso è \emph{assente}, ed è la condizione che il braccio batch
implementa.

\paragraph{Il simulatore CARLA.}
CARLA~\cite{dosovitskiy2017carla} è un simulatore open-source di guida
urbana su Unreal Engine~4, distribuito come server che ospita mondo e
fisica e pilotato da client via API Python. Conta qui per tre ragioni
(versione 0.9.16 ovunque): la \emph{modalità sincrona a passo fisso}
rende la raccolta di esperienza indipendente dal rendering; i pedoni
sono attori \emph{walker} ad azionamento diretto, il che rende possibile
apprenderne le policy; e le città fornite differiscono \emph{in natura}
e non solo in dimensione --- l'addestramento usa \emph{Town03}, ``a
larger, urban map with a roundabout and large junctions'', il test di
generalizzazione \emph{Town05}, una ``squared-grid town with cross
junctions and a bridge'' a più corsie per direzione~\cite{carladocs}.
Anche in modalità sincrona e con tutti i seed fissati, infine, CARLA
sotto carico multi-agente non è deterministico da run a run: una
variabilità irriducibile quantificata e usata come metro per ogni
effetto riportato (\cref{sec:gate}).

\section{Approccio, contributi e struttura}
\label{sec:approccio}

Tutti gli esperimenti girano su Town03, con Town05 riservata al solo
test di generalizzazione. Sei agenti apprendono simultaneamente in un
unico mondo condiviso --- tre veicoli e tre pedoni, ciascuno con un
percorso proprio a ogni episodio --- un design che chiamiamo
\emph{copertura parallela dei task} e che è la cornice fissa di entrambi
i bracci sperimentali, non una variabile. Le policy sono addestrate con
MAPPO in regime CTDE, con una policy condivisa per classe e un critic
centralizzato che vede le osservazioni di tutti gli agenti solo durante
l'addestramento; la difficoltà è la lunghezza target del percorso
(30/60/\SI{100}{m}) e il successo è il completamento rigoroso del
percorso assegnato.

I due regimi condividono ogni componente tranne la regola di selezione
del livello: il braccio \emph{curriculum} parte dal solo easy e sblocca
i livelli successivi attraverso un gate di competenza a finestra, il
braccio \emph{batch} alterna uniformemente tutti i livelli fin dal primo
episodio. La campagna finale è un fattoriale $2 \times 2$ --- \{critic
MLP, critic GNN\} $\times$ \{curriculum, batch\} --- di quattro run da
tre milioni di passi a seed singolo appaiato, ciascuna seguita da una
valutazione di 400 episodi che include Town05. Con una run per cella
ogni contrasto è un caso di studio appaiato, letto contro una banda di
varianza misurata su repliche a seed identico; il contrasto primario è
però osservato \emph{due volte}, una per architettura del critic.
Anticipando il \cref{ch:risultati}: la clausola di efficienza di H1 e
quella di robustezza di H2 falliscono entrambe, mentre H3 regge in forma
forte.

I contributi sono quattro: \textbf{(i)} un framework di confronto
controllato a parità di budget, con run appaiate per seed e regole di
misura rigorose per agente; \textbf{(ii)} una metodologia di sviluppo
guidata da gate con un registro onesto di modifiche promosse, rigettate
e consapevolmente mantenute, in cui i difetti di validità dell'ambiente
--- incluso uno che ha spostato la metrica principale di $\approx 50$
punti percentuali senza alcun cambiamento di policy --- sono documentati
con le loro firme di rilevamento; \textbf{(iii)} una caratterizzazione
empirica di che cosa fa l'ordinamento dei task, ossia nessun effetto di
efficienza iniziale ma selezione tra equilibri di guida
qualitativamente diversi; e \textbf{(iv)} l'evidenza che gli effetti di
regime sono condizionati dalla stima del valore.

La tesi prosegue in quattro blocchi. I
\cref{ch:piattaforma,ch:metodo,ch:risultati} descrivono il progetto:
rispettivamente la piattaforma sperimentale, il metodo di misura con il
design della campagna, e i risultati con la loro interpretazione. Il
\cref{ch:evoluzioni} raccoglie le possibili evoluzioni del lavoro e il
\cref{ch:conclusioni} conclude.
